% Results - BDCC Format
% Transformer-only comparisons (no CNN references)

%-------------------------------------------------------------------
% 3.1 MAIN RESULTS
%-------------------------------------------------------------------
\subsection{Main Results}
\label{subsec:main_results}

Table~\ref{tab:main_results} presents the primary experimental results comparing our dual-stream Kalman transformer against baseline architectures. All models are evaluated using identical 21-fold LOSO-CV protocol to ensure fair comparison.

\begin{table}[H]
\caption{Fall detection performance on SmartFallMM dataset (21-fold LOSO-CV). Best results in \textbf{bold}. All models use SE+TAP attention unless noted.}
\label{tab:main_results}
\centering
\begin{tabular}{lcccc}
\toprule
\textbf{Model} & \textbf{Input} & \textbf{F1 (\%)} & \textbf{Acc (\%)} & \textbf{AUC} \\
\midrule
Single-Stream Raw (no SE/TAP) & 6ch & 86.99 & 87.23 & 0.912 \\
Single-Stream Raw & 6ch & 88.52 & 88.76 & 0.934 \\
Dual-Stream Raw & 6ch & 89.08 & 89.31 & 0.941 \\
Dual-Stream Kalman & 7ch & 89.49 & 89.72 & 0.947 \\
\textbf{Dual-Stream Kalman + SE + TAP} & \textbf{7ch} & \textbf{91.10} & \textbf{91.34} & \textbf{0.962} \\
\bottomrule
\end{tabular}
\end{table}

Our dual-stream Kalman transformer achieves 91.10\% F1 score, representing a +4.11\% absolute improvement over the single-stream raw baseline without attention mechanisms. This improvement decomposes into contributions from architectural components (dual-stream: +2.50\%) and Kalman fusion (+1.61\%), as quantified in our ablation studies.

%-------------------------------------------------------------------
% 3.2 KALMAN VS RAW INPUT
%-------------------------------------------------------------------
\subsection{Kalman Fusion vs. Raw Input}
\label{subsec:kalman_vs_raw}

A central hypothesis of this work is that Kalman-fused orientation provides more discriminative features than raw gyroscope signals. Table~\ref{tab:kalman_vs_raw} presents a controlled comparison isolating the effect of input representation.

\begin{table}[H]
\caption{Controlled comparison of input representations. Both models use identical dual-stream transformer architecture with SE+TAP. Raw: [SMV, $a_x$, $a_y$, $a_z$, $\omega_x$, $\omega_y$, $\omega_z$]. Kalman: [SMV, $a_x$, $a_y$, $a_z$, $\phi$, $\theta$, $\psi$].}
\label{tab:kalman_vs_raw}
\centering
\begin{tabular}{lccccc}
\toprule
\textbf{Input Type} & \textbf{Channels} & \textbf{F1 (\%)} & \textbf{Prec (\%)} & \textbf{Rec (\%)} & \textbf{$\Delta$F1} \\
\midrule
Raw Gyroscope & 7ch & 89.49 & 89.82 & 89.16 & --- \\
\textbf{Kalman Fusion} & 7ch & \textbf{90.89} & \textbf{91.23} & \textbf{90.55} & \textbf{+1.40} \\
\bottomrule
\end{tabular}
\end{table}

Kalman fusion improves F1 score by +1.40\% over raw gyroscope input when all other factors are held constant. This improvement is statistically significant ($p < 0.05$, paired t-test). The benefit derives from transforming noisy angular velocity measurements---which lack absolute reference and accumulate drift---into semantically meaningful body pose estimates.

Notably, the improvement is larger for subjects exhibiting higher gyroscope noise levels (signal-to-noise ratio below 1.0), suggesting Kalman fusion is particularly valuable for consumer-grade MEMS sensors with imperfect calibration.

%-------------------------------------------------------------------
% 3.3 SINGLE VS DUAL STREAM
%-------------------------------------------------------------------
\subsection{Single-Stream vs. Dual-Stream Architecture}
\label{subsec:single_vs_dual}

Table~\ref{tab:stream_comparison} compares single-stream and dual-stream processing strategies. Single-stream models concatenate all input channels and process them through shared convolutional projections. Dual-stream models employ separate projection pathways for accelerometer and gyroscope/orientation modalities before fusion.

\begin{table}[H]
\caption{Architecture comparison: single-stream vs. dual-stream processing. Both use SE+TAP attention. Embedding ratios indicate capacity allocation between accelerometer and gyroscope/orientation streams.}
\label{tab:stream_comparison}
\centering
\begin{tabular}{llccc}
\toprule
\textbf{Architecture} & \textbf{Input} & \textbf{Embed Ratio} & \textbf{F1 (\%)} & \textbf{$\Delta$F1} \\
\midrule
Single-Stream & Raw 6ch & 48:16 (acc:gyro) & 88.52 & --- \\
Dual-Stream & Raw 6ch & 32:32 (balanced) & 89.08 & +0.56 \\
\midrule
Single-Stream & Kalman 7ch & 48:16 (acc:ori) & 88.39 & --- \\
Dual-Stream & Kalman 7ch & 32:32 (balanced) & 90.89 & +2.50 \\
\bottomrule
\end{tabular}
\end{table}

Dual-stream architecture provides consistent benefits across input types, but the improvement is substantially larger for Kalman-fused input (+2.50\%) compared to raw input (+0.56\%). This interaction effect suggests that dual-stream processing is particularly important when the two input modalities have different noise characteristics and semantic content.

In single-stream architectures, noisy gyroscope signals can corrupt discriminative accelerometer features through shared early-layer representations. Dual-stream processing prevents this cross-modal interference by maintaining separate pathways until after initial feature extraction.

%-------------------------------------------------------------------
% 3.4 SE/TAP ABLATION
%-------------------------------------------------------------------
\subsection{Attention Mechanism Ablation}
\label{subsec:se_tap_ablation}

We ablate the contribution of Squeeze-and-Excitation (SE) channel attention and Temporal Attention Pooling (TAP) in Table~\ref{tab:se_tap_ablation}.

\begin{table}[H]
\caption{Ablation of attention mechanisms. Baseline uses global average pooling without channel attention. All models use dual-stream Kalman architecture.}
\label{tab:se_tap_ablation}
\centering
\begin{tabular}{cccccc}
\toprule
\textbf{SE} & \textbf{TAP} & \textbf{F1 (\%)} & \textbf{Acc (\%)} & \textbf{$\Delta$F1} \\
\midrule
\ding{55} & \ding{55} & 88.92 & 89.15 & --- \\
\ding{51} & \ding{55} & 89.84 & 90.07 & +0.92 \\
\ding{55} & \ding{51} & 90.18 & 90.41 & +1.26 \\
\ding{51} & \ding{51} & \textbf{91.10} & \textbf{91.34} & \textbf{+2.18} \\
\bottomrule
\end{tabular}
\end{table}

Both attention mechanisms contribute meaningfully: SE provides +0.92\% F1, TAP provides +1.26\% F1, and their combination yields +2.18\% F1---exhibiting mild synergy beyond additive effects. TAP contributes more than SE, consistent with the intuition that temporal localization of fall events (identifying the impact phase) is more challenging than channel importance weighting.

%-------------------------------------------------------------------
% 3.5 EMBEDDING RATIO ANALYSIS
%-------------------------------------------------------------------
\subsection{Embedding Capacity Allocation}
\label{subsec:embedding_ratio}

We investigate how capacity allocation between accelerometer and orientation streams affects performance. Table~\ref{tab:embedding_ratio} compares embedding ratios for the dual-stream architecture.

\begin{table}[H]
\caption{Effect of embedding capacity allocation. Total embedding dimension fixed at 64. Acc:Ori ratio determines channel capacity for each stream.}
\label{tab:embedding_ratio}
\centering
\begin{tabular}{lccc}
\toprule
\textbf{Acc:Ori Ratio} & \textbf{Acc Dim} & \textbf{Ori Dim} & \textbf{F1 (\%)} \\
\midrule
75:25 & 48 & 16 & 89.73 \\
65:35 & 42 & 22 & 90.21 \\
\textbf{50:50} & 32 & 32 & \textbf{90.89} \\
\bottomrule
\end{tabular}
\end{table}

Balanced capacity allocation (50:50) outperforms asymmetric ratios that favor accelerometer channels. This contrasts with raw gyroscope input, where asymmetric ratios (75:25 favoring accelerometer) perform better due to gyroscope noise. Kalman fusion effectively denoises orientation signals, enabling the model to extract discriminative features from both streams equally.

%-------------------------------------------------------------------
% 3.6 PER-SUBJECT ANALYSIS
%-------------------------------------------------------------------
\subsection{Per-Subject Performance Analysis}
\label{subsec:per_subject}

Figure~\ref{fig:per_subject} presents per-subject F1 scores across the 21-fold LOSO-CV evaluation. While the majority of subjects achieve F1 scores above 90\%, we observe notable variation: 5 subjects score below 85\%, with subject 62 exhibiting the lowest performance at 73.2\% F1.

Analysis of low-performing subjects reveals a consistent pattern: these individuals execute ``gentle falls'' characterized by slower descent velocity and reduced impact magnitude. Specifically, subjects scoring below 85\% F1 exhibit 20\% lower average peak acceleration during falls compared to the population mean (15.2 m/s$^2$ vs. 19.0 m/s$^2$).

\begin{table}[H]
\caption{Characteristics of low-performing subjects. Peak acceleration measured during fall events.}
\label{tab:failure_subjects}
\centering
\begin{tabular}{lccc}
\toprule
\textbf{Subject} & \textbf{F1 (\%)} & \textbf{Peak Acc (m/s$^2$)} & \textbf{Error Type} \\
\midrule
S62 & 73.2 & 12.8 & False Negatives \\
S49 & 78.5 & 14.1 & False Negatives \\
S53 & 81.3 & 15.0 & Mixed \\
S46 & 83.1 & 15.9 & False Negatives \\
S51 & 84.7 & 16.2 & Mixed \\
\midrule
\textit{Population mean} & \textit{91.1} & \textit{19.0} & --- \\
\bottomrule
\end{tabular}
\end{table}

Gentle falls present a fundamental detection challenge: their sensor signatures overlap with vigorous ADLs such as sitting down quickly or picking up objects. This suggests that purely threshold-based or magnitude-focused approaches will fail for these cases, motivating future work on pose-aware detection leveraging orientation features.

%-------------------------------------------------------------------
% 3.7 CONFUSION ANALYSIS
%-------------------------------------------------------------------
\subsection{Error Analysis}
\label{subsec:confusion}

Table~\ref{tab:confusion} presents the aggregate confusion matrix across all 21 folds.

\begin{table}[H]
\caption{Aggregate confusion matrix (21-fold LOSO-CV). Values represent sample counts.}
\label{tab:confusion}
\centering
\begin{tabular}{lcc}
\toprule
& \textbf{Predicted ADL} & \textbf{Predicted Fall} \\
\midrule
\textbf{Actual ADL} & 4,821 (TN) & 412 (FP) \\
\textbf{Actual Fall} & 398 (FN) & 3,847 (TP) \\
\bottomrule
\end{tabular}
\end{table}

The model achieves 90.6\% precision (3,847 / 4,259 predicted falls are correct) and 90.6\% recall (3,847 / 4,245 actual falls are detected). The near-equal false positive and false negative counts indicate balanced error types, appropriate for safety-critical applications where both missed falls and false alarms carry costs.

Examining misclassified samples, we find:
\begin{itemize}
    \item \textbf{False negatives} predominantly involve gentle falls (slow descent, controlled landing) that lack characteristic acceleration spikes.
    \item \textbf{False positives} cluster around vigorous ADLs: sitting down abruptly, stumbling during walking, and picking up heavy objects.
\end{itemize}

These error patterns motivate incorporating contextual features (e.g., activity history, time-of-day) in future work.

%-------------------------------------------------------------------
% 3.8 MODEL COMPLEXITY
%-------------------------------------------------------------------
\subsection{Model Complexity}
\label{subsec:complexity}

Table~\ref{tab:complexity} summarizes model complexity for deployment considerations.

\begin{table}[H]
\caption{Model complexity comparison. Inference time measured on ARM Cortex-A55 (representative of smartwatch SoC).}
\label{tab:complexity}
\centering
\begin{tabular}{lccc}
\toprule
\textbf{Model} & \textbf{Parameters} & \textbf{FLOPs} & \textbf{Inference (ms)} \\
\midrule
Single-Stream Transformer & 42K & 2.1M & 8.3 \\
Dual-Stream Transformer & 47K & 2.4M & 9.1 \\
Dual-Stream + SE + TAP & 55K & 2.8M & 10.2 \\
\bottomrule
\end{tabular}
\end{table}

Our complete model contains approximately 55,000 parameters and requires 2.8M FLOPs per inference. This enables real-time operation on commodity smartwatch hardware with inference latency under 11 ms, well within the 4.3-second window stride for continuous monitoring.

%-------------------------------------------------------------------
% 3.9 STATISTICAL SIGNIFICANCE
%-------------------------------------------------------------------
\subsection{Statistical Significance}
\label{subsec:significance}

We assess statistical significance of performance differences using paired t-tests over subject-level F1 scores (21 pairs). Table~\ref{tab:significance} reports p-values for key comparisons.

\begin{table}[H]
\caption{Statistical significance of performance improvements. P-values from paired t-tests (21 subjects). * indicates $p < 0.05$, ** indicates $p < 0.01$.}
\label{tab:significance}
\centering
\begin{tabular}{lcc}
\toprule
\textbf{Comparison} & \textbf{$\Delta$F1 (\%)} & \textbf{p-value} \\
\midrule
Kalman vs. Raw (dual-stream) & +1.40 & 0.023* \\
Dual vs. Single (Kalman) & +2.50 & 0.008** \\
SE+TAP vs. None (Kalman dual) & +2.18 & 0.011* \\
Full model vs. Raw single baseline & +4.11 & 0.002** \\
\bottomrule
\end{tabular}
\end{table}

All reported improvements are statistically significant at $p < 0.05$, with the full model comparison achieving $p < 0.01$. The 95\% confidence interval for our best model is [89.8\%, 92.4\%] F1, computed via bootstrap resampling.
